\documentclass[a4paper]{article}
\usepackage{amsfonts}
\usepackage{amsmath}
\usepackage{amssymb}
\usepackage{graphicx}  

\newcommand{\degree}{^{\circ}}
\newcommand{\cis}{\operatorname{cis}}
\newcommand{\Bgsin}{\operatorname{Bgsin}}
\newcommand{\Bgcos}{\operatorname{Bgcos}}
\newcommand{\Bgtan}{\operatorname{Bgtan}}

\begin{document}
  
\noindent \large Auteur: Sadia Nazary  \\
\noindent \large Studierichting: Informatica\\
\noindent \large Oefeningenreeks 1E nr. 2\\

\medskip

\normalsize

We weten dat $R(z)$ van de derde graad moet zijn, dus: \\

$R(z)=az^3+bz^2+cz+d$ \\

\medskip

We schrijven: \\

$z^6+3z^5+2z^4+5z^3-6z^2-8z+3=(z^4-1)Q(z)+R(z)$ \\

\medskip

De nulpunten van $z^4-1=0$ zijn: \\

$z=1,-1,i,-i$ \\

\medskip

Invullen van $z=1$: \\

$a+b+c+d=1+3+2+5-6-8+3$ \\

$a+b+c+d=0$ \\

\medskip

Invullen van $z=-1$: \\

$-a+b-c+d=1-3+2-5-6+8+3$ \\

$-a+b-c+d=0$ \\

\medskip

Invullen van $z=i$: \\

$-ai-b+ci+d=10-10i$ \\

\medskip

Invullen van $z=-i$: \\

$ai-b-ci+d=10+10i$ \\

\medskip

We bekomen het stelsel: \\

$a+b+c+d=0$ \\

$-a+b-c+d=0$ \\

$(c-a)i+d-b=10-10i$ \\

$(a-c)i+d-b=10+10i$ \\

\medskip

Uit de eerste twee vergelijkingen: \\

$2a+2c=0$ \\

$a=-c$ \\

\medskip

Verder: \\

$c-a=-10$ \\

$2c=-10$ \\

$c=-5$ \\

$a=5$ \\

\medskip

$b=-d$ \\

$d-b=10$ \\

$2d=10$ \\

$d=5$ \\

$b=-5$ \\

\medskip

Dus: \\

$a=5$ \\

$b=-5$ \\

$c=-5$ \\

$d=5$ \\

\medskip

$R(z)=5z^3-5z^2-5z+5$ \\

\begin{thebibliography}{999}
\end{thebibliography}

\end{document}